% EXTRACTED from v1.1 manuscript (tag v1.1-verified): Part II, ch. "Natural
% Deduction and Sequents" sec. "Classical sequent calculus" (LK_CPL) and
% Part III, ch. "First-Order Deduction" sec. "First-order sequent rules"
% (LK_FOL quantifier rules).  Per the Phase-3 nomenclature ruling, sheet:lk is
% this two-sided, primitive-negation, contexted-initial-sequent presentation.
% Mix (multicut) is a Cut-equivalent derived convenience for this canonical,
% primitive-negation presentation.  The Phase-4 Part-7 cut-elimination sketch
% is NOT a theorem about this sheet: it is scoped to sheet:lk-defneg and
% sheet:lj, whose logical rule cases are the ones treated by v1.1 Part VIII.
\calculussheet{sheet:lk}{Classical Sequent Calculus $\mathsf{LK}$ (contexted initial sequents)}{Language: $\mathsf{Fm}(\Lang)$, first-order with primitives $\bot,\neg,\land,\lor,\to,\forall,\exists$; sequent shape: two-sided sequents $\Gamma\Rightarrow\Delta$ with finite multisets on both sides (exchange built into the multiset convention), read as $\bigwedge\Gamma\to\bigvee\Delta$ with empty antecedent $\top$ and empty succedent $\bot$; initial sequents are contexted ($\Gamma,\varphi\Rightarrow\varphi,\Delta$ and $\Gamma,\bot\Rightarrow\Delta$); negation: primitive, with rules $\neg L$/$\neg R$; equality: none (sheet:lk-eq adds it as a distinct calculus).}
\hbrule{rule:lk-id}{id}{}{\Gamma,\varphi\Rightarrow\varphi,\Delta}{\none}
\hbrule{rule:lk-bot-l}{$\bot L$}{}{\Gamma,\bot\Rightarrow\Delta}{\none}
\hbrule{rule:lk-wl}{$W_L$}{\Gamma\Rightarrow\Delta}{\Gamma,\varphi\Rightarrow\Delta}{\none}
\hbrule{rule:lk-wr}{$W_R$}{\Gamma\Rightarrow\Delta}{\Gamma\Rightarrow\Delta,\varphi}{\none}
\hbrule{rule:lk-cl}{$C_L$}{\Gamma,\varphi,\varphi\Rightarrow\Delta}{\Gamma,\varphi\Rightarrow\Delta}{\none}
\hbrule{rule:lk-cr}{$C_R$}{\Gamma\Rightarrow\Delta,\varphi,\varphi}{\Gamma\Rightarrow\Delta,\varphi}{\none}
\hbrule{rule:lk-neg-l}{$\neg L$}{\Gamma\Rightarrow\Delta,\varphi}{\Gamma,\neg\varphi\Rightarrow\Delta}{\none}
\hbrule{rule:lk-neg-r}{$\neg R$}{\Gamma,\varphi\Rightarrow\Delta}{\Gamma\Rightarrow\Delta,\neg\varphi}{\none}
\hbrule{rule:lk-and-l}{$\land L$}{\Gamma,\varphi,\psi\Rightarrow\Delta}{\Gamma,\varphi\land\psi\Rightarrow\Delta}{\none}
\hbrule{rule:lk-and-r}{$\land R$}{\Gamma\Rightarrow\Delta,\varphi\qquad\Gamma\Rightarrow\Delta,\psi}{\Gamma\Rightarrow\Delta,\varphi\land\psi}{\none}
\hbrule{rule:lk-or-l}{$\lor L$}{\Gamma,\varphi\Rightarrow\Delta\qquad\Gamma,\psi\Rightarrow\Delta}{\Gamma,\varphi\lor\psi\Rightarrow\Delta}{\none}
\hbrule{rule:lk-or-r}{$\lor R$}{\Gamma\Rightarrow\Delta,\varphi,\psi}{\Gamma\Rightarrow\Delta,\varphi\lor\psi}{\none}
\hbrule{rule:lk-imp-l}{$\to L$}{\Gamma\Rightarrow\Delta,\varphi\qquad\Gamma,\psi\Rightarrow\Delta}{\Gamma,\varphi\to\psi\Rightarrow\Delta}{\none}
\hbrule{rule:lk-imp-r}{$\to R$}{\Gamma,\varphi\Rightarrow\Delta,\psi}{\Gamma\Rightarrow\Delta,\varphi\to\psi}{\none}
\hbrule{rule:lk-cut}{Cut}{\Gamma\Rightarrow\Delta,\varphi\qquad\Sigma,\varphi\Rightarrow\Pi}{\Gamma,\Sigma\Rightarrow\Delta,\Pi}{\none}
\hbrule{rule:lk-forall-l}{$\forall L$}{\Gamma,\varphi[t/x]\Rightarrow\Delta}{\Gamma,\forall x\,\varphi\Rightarrow\Delta}{substitution: $[t/x]$ is the total capture-avoiding substitution fixed in the syntax chapter, so no separate substitutability proviso is imposed}
\hbrule{rule:lk-forall-r}{$\forall R$}{\Gamma\Rightarrow\Delta,\varphi[y/x]}{\Gamma\Rightarrow\Delta,\forall x\,\varphi}{eigenvariable freshness: $y$ occurs free in no formula of $\Gamma,\Delta,\forall x\,\varphi$; substitution total and capture-avoiding}
\hbrule{rule:lk-exists-l}{$\exists L$}{\Gamma,\varphi[y/x]\Rightarrow\Delta}{\Gamma,\exists x\,\varphi\Rightarrow\Delta}{eigenvariable freshness: $y$ occurs free in no formula of $\Gamma,\Delta,\exists x\,\varphi$; substitution total and capture-avoiding}
\hbrule{rule:lk-exists-r}{$\exists R$}{\Gamma\Rightarrow\Delta,\varphi[t/x]}{\Gamma\Rightarrow\Delta,\exists x\,\varphi}{substitution: $[t/x]$ is the total capture-avoiding substitution, as in $\forall L$}
\hbrule{rule:lk-mix}{Mix}{\Gamma\Rightarrow\Delta,\varphi^{m}\qquad\varphi^{n},\Sigma\Rightarrow\Pi}{\Gamma,\Sigma\Rightarrow\Delta,\Pi}{$m,n\geq 1$; $\varphi^{k}$ abbreviates $k$ comma-separated occurrences of the mix formula $\varphi$; $\Delta$ and $\Sigma$ may contain further occurrences of $\varphi$; this rule is equipollent with primitive Cut by contraction, but the Part-7 cut-elimination sketch does not cover this primitive-negation presentation}

The defined-negation translation is
Definition~\ref{def:p07-defined-negation-translation}; its derivability
comparison with Sheet~\ref{sheet:lk-defneg} is
Theorem~\ref{thm:p07-lk-presentation-equivalence}, with the explicitly
sketched proof in Annex~\ref{annex:p7.calculi.2}.  This comparison does not
identify proof trees or transfer a cut-elimination proof automatically.
